Nachrichten erreichen ihr Publikum heute zu einem großen Teil nicht mehr über die Startseite einer Zeitung, sondern über soziale Plattformen, Suchmaschinen und Smartphone-Apps.\kihinweis{}
Der Reuters Institute Digital News Report 2025 zeigt, dass der Anteil der Befragten, die pro Woche Nachrichten in Form von Social-Media-Videos nutzen, über alle untersuchten Märkte von 52\,\% im Jahr 2020 auf 65\,\% im Jahr 2025 gestiegen ist \cite{newman2025dnr}.
Bei den 18- bis 24-Jährigen geben 44\,\% an, dass soziale Netzwerke und Videoplattformen ihre wichtigste Nachrichtenquelle sind \cite{newman2025dnr}.
Plattformen wie TikTok oder Instagram haben dabei ein eigenes Format geprägt: kurze Inhalte im Hochformat, die einzeln den ganzen Bildschirm füllen und mit einer Wischbewegung gewechselt werden.

Klassische Nachrichtenangebote passen nur bedingt zu diesem Nutzungsverhalten.
Ein Artikel auf einer Nachrichtenseite besteht aus Überschrift, Vorspann und mehreren Absätzen Fließtext und ist für das Lesen am Stück gedacht.
Um denselben Inhalt in einem Short-Form-Feed anzubieten, braucht es zusätzlich eine kurze Fassung, ein passendes Bild im Hochformat und eine Einordnung in Themen.
Werden diese Inhalte von Hand erstellt, entsteht für jeden Artikel zusätzlicher Aufwand in der Redaktion.
Sprachmodelle können diese Aufgabe inzwischen in einer Qualität übernehmen, die in Studien mit von Menschen geschriebenen Zusammenfassungen vergleichbar ist \cite{zhang2024benchmarking}.
Gleichzeitig bleiben sie fehleranfällig und erfinden mitunter Inhalte, die nicht im Ausgangstext stehen \cite{maynez2020faithfulness}.

Aus dieser Ausgangslage ist gemeinsam mit dem Projektpartner Werbeberg die vorliegende Diplomarbeit entstanden.
Ziel ist ein System, das aus bestehenden Nachrichtenartikeln automatisch einen Teaser, Kernpunkte, Themen-Tags und ein Titelbild erzeugt und diese Inhalte als vertikalen Feed ausliefert.
Der Feed soll nicht als eigenständige App enden, sondern als Komponente in bestehende Nuxt-Websites eingebaut werden können.
Die generierten Inhalte sollen nachvollziehbar und reproduzierbar bleiben und von einer Redaktion geprüft und korrigiert werden können.

Die Arbeit gliedert sich wie folgt.
Kapitel~\ref{chapter:related} ordnet das Projekt in bestehende Angebote und in die Forschung zu automatischer Zusammenfassung und Personalisierung ein.
Kapitel~\ref{chapter:requirements} beschreibt den Projektauftrag und die daraus abgeleiteten Anforderungen.
Kapitel~\ref{chapter:architecture} gibt einen Überblick über die Bausteine des Systems und deren Zusammenspiel, bevor die folgenden Kapitel auf die Umsetzung im Detail eingehen.
